% Created 2026-02-24 Tue 22:57
% Intended LaTeX compiler: pdflatex
\documentclass[11pt]{article}
\usepackage[utf8]{inputenc}
\usepackage[T1]{fontenc}
\usepackage{graphicx}
\usepackage{longtable}
\usepackage{wrapfig}
\usepackage{rotating}
\usepackage[normalem]{ulem}
\usepackage{amsmath}
\usepackage{amssymb}
\usepackage{capt-of}
\usepackage{hyperref}
\author{The Prologos Project — from Zachary Larson}
\date{2026-02-21}
\title{Prologos Language Grammar\\\medskip
\large A Literate Specification of the Surface Syntax}
\hypersetup{
 pdfauthor={The Prologos Project — from Zachary Larson},
 pdftitle={Prologos Language Grammar},
 pdfkeywords={},
 pdfsubject={},
 pdfcreator={Emacs 30.2 (Org mode 9.7.39)}, 
 pdflang={English}}
\begin{document}

\maketitle
\setcounter{tocdepth}{3}
\tableofcontents

\section{Introduction}
\label{sec:org9942054}

Prologos is a functional-logic language unifying dependent types, session types,
linear types (QTT), logic programming, and propagators. This document describes
the \emph{surface syntax} of the language as written in \texttt{.prologos} files.

Prologos has two syntactic modes:

\begin{enumerate}
\item \textbf{Whitespace-sensitive mode} (\texttt{.prologos} files): Indentation-based structure
with \texttt{[]} for grouping and minimal punctuation. This is the primary mode.
\item \textbf{S-expression mode}: Parenthesized fallback for use in macros, tests, and
when embedding in Racket. Every WS-mode form has a canonical sexp
representation.
\end{enumerate}

The grammar is organized bottom-up: lexical elements, then types, then
expressions, then declarations.

\begin{quote}
\textbf{Design Principle}: Prologos is \emph{homoiconic} --- code and data share the same
representation. All syntactic sugar desugars to s-expressions. Macros operate on
the post-parse representation, making code-as-data a first-class concept.
\end{quote}
\section{Lexical Grammar}
\label{sec:org90a0a46}

\subsection{Comments}
\label{sec:org903ac13}

Line comments begin with \texttt{;} and extend to end of line. There are no block
comments.

\begin{verbatim}
;; This is a comment
def x : Nat zero  ; inline comment
\end{verbatim}
\subsection{Identifiers}
\label{sec:org425539b}

Identifiers may contain letters, digits, hyphens, question marks, exclamation
marks, and primes. They must not start with a digit.

\begin{verbatim}
identifier  ::= ident-start { ident-char }
ident-start ::= letter | '_' | '+' | '*' | '/' | '!' | '?' | '<' | '>' | '='
ident-char  ::= ident-start | digit | '-' | "'"
\end{verbatim}

Naming conventions:
\begin{itemize}
\item Predicates use \texttt{?} suffix: \texttt{zero?}, \texttt{empty?}, \texttt{member?}
\item Destructive operations use \texttt{!} suffix: \texttt{persist!}, \texttt{tvec-push!}
\item Type names are capitalized: \texttt{Nat}, \texttt{List}, \texttt{Option}
\item Value names are lowercase with hyphens: \texttt{sort-on}, \texttt{find-index}
\end{itemize}

Qualified names use \texttt{::} separator:

\begin{verbatim}
list::map       ;; qualified reference to list module's map
nat::add        ;; qualified reference to nat module's add
opt::unwrap-or  ;; qualified reference to option module's unwrap-or
\end{verbatim}
\subsection{Numeric Literals}
\label{sec:org9092880}

Prologos has four numeric literal forms, plus a decimal variant for approximate literals:

\begin{verbatim}
nat-literal     ::= digit+ 'N'                (* Natural: 0N, 42N, 100N *)
int-literal     ::= ['-'] digit+              (* Integer: 42, -7, 0     *)
rat-literal     ::= ['-'] digit+ '/' digit+   (* Rational: 3/7, -1/3   *)
decimal-literal ::= ['-'] digit+ '.' digit+   (* Decimal: 3.14, 0.5    *)
approx-literal  ::= '~' (int | rat | decimal) (* Posit approx: ~42, ~3/7, ~3.14 *)
\end{verbatim}

\begin{verbatim}
42N          ;; Nat literal (Church-encoded natural — type infrastructure only)
42           ;; Int literal (arbitrary-precision integer)
3/7          ;; Rat literal (exact rational)
3.14         ;; Posit32 literal (approximate — same as ~3.14)
~3/7         ;; approximate Posit literal (from fraction)
~3.14        ;; approximate Posit literal (from decimal, stored as exact rational 157/50)
\end{verbatim}

Bare decimal literals (\texttt{3.14}, \texttt{0.5}) produce Posit32 values, the same as their tilde-prefixed
equivalents (\texttt{\textasciitilde{}3.14}, \texttt{\textasciitilde{}0.5}). Internally, the decimal is stored as an exact rational (\texttt{157/50})
and encoded to the nearest Posit32 bit pattern.

Nat is intended for type-level infrastructure (indices, lengths, proofs), not general computation.
\subsection{String Literals}
\label{sec:orgdef83ec}

Double-quoted with standard escape sequences (\texttt{\textbackslash{}n}, \texttt{\textbackslash{}t}, \texttt{\textbackslash{}\textbackslash{}}, \texttt{\textbackslash{}"}).

\begin{verbatim}
"hello world"
"line one\nline two"
\end{verbatim}
\subsection{Keyword Literals}
\label{sec:org38b1556}

Keywords start with \texttt{:} and are used as map keys and enum-like values:

\begin{verbatim}
:name
:age
:hello
\end{verbatim}
\subsection{Boolean Literals}
\label{sec:org723d232}

\begin{verbatim}
true
false
\end{verbatim}
\subsection{Special Tokens}
\label{sec:org9750138}

\begin{center}
\begin{tabular}{ll}
Token & Meaning\\
\hline
\texttt{\_} & Wildcard / type hole (inferred)\\
\texttt{\_1} \texttt{\_2} & Numbered placeholders for partial application\\
\texttt{??} & Typed hole (interactive development)\\
\texttt{??name} & Named typed hole\\
\texttt{:0} & Erased multiplicity (use 0 times)\\
\texttt{:1} & Linear multiplicity (use exactly 1 time)\\
\texttt{:w} & Unrestricted multiplicity (use any times)\\
\texttt{zero} & Nat zero constructor\\
\texttt{unit} & Unit value\\
\texttt{refl} & Equality reflexivity proof\\
\end{tabular}
\end{center}
\subsection{Bracket Types}
\label{sec:org80089dc}

Prologos uses four bracket types, each with distinct semantics:

\begin{center}
\begin{tabular}{ll}
Brackets & Purpose\\
\hline
\texttt{[...]} & Primary grouping: function application, params\\
\texttt{(...)} & Special forms: \texttt{(fn ...)}, \texttt{(match ...)}, types\\
\texttt{\{...\}} & Implicit type parameters, map literals\\
\texttt{<...>} & Dependent types, return type annotations\\
\end{tabular}
\end{center}

Inside any bracket pair, newlines are treated as whitespace (indentation is not
significant).
\subsection{Collection Literal Prefixes}
\label{sec:orgd967663}

\begin{center}
\begin{tabular}{lll}
Syntax & Type & Example\\
\hline
\texttt{'[...]} & \texttt{List} & \texttt{'[1N 2N 3N]}\\
\texttt{@[...]} & \texttt{PVec} & \texttt{@[1 2 3]}\\
\texttt{\textasciitilde{}[...]} & \texttt{LSeq} & \texttt{\textasciitilde{}[1 2 3]}\\
\texttt{\#\{...\}} & \texttt{Set} & \texttt{\#\{1 2 3\}}\\
\texttt{\{k v\}} & \texttt{Map} & \texttt{\{:name "Ada"\}}\\
\end{tabular}
\end{center}
\section{Type Expressions}
\label{sec:orgeb03ab5}

In a dependently-typed language, types and terms share the same expression
syntax. This section highlights type-specific forms.
\subsection{Base Types}
\label{sec:org1a9b3ef}

\begin{verbatim}
base-type ::= 'Type'     (* universe of types, level inferred *)
            | 'Nat'      (* natural numbers                    *)
            | 'Int'      (* arbitrary-precision integers        *)
            | 'Rat'      (* exact rationals                     *)
            | 'Bool'     (* booleans                            *)
            | 'Unit'     (* unit type                           *)
            | 'Symbol'   (* symbols for homoiconicity           *)
            | 'Keyword'  (* keyword type for map keys           *)
            | 'Datum'    (* code-as-data algebraic type         *)
\end{verbatim}

\begin{verbatim}
def x : Nat zero
def b : Bool true
def u : Unit unit
\end{verbatim}
\subsection{Parameterized Types}
\label{sec:orgcebf12b}

Types are applied to arguments via juxtaposition in brackets:

\begin{verbatim}
List Nat            ;; list of natural numbers
Option Bool         ;; optional boolean
Map Keyword Int     ;; keyword-to-int map
Vec Nat 3N          ;; length-3 vector of Nats (dependent!)
Set Int             ;; set of ints
PVec Nat            ;; persistent vector of Nats
\end{verbatim}
\subsection{Function Types (Arrows)}
\label{sec:org565b9c6}

Arrows are right-associative infix operators:

\begin{verbatim}
arrow-type ::= type '->' type    (* unrestricted function *)
             | type '-0>' type   (* erased function       *)
             | type '-1>' type   (* linear function       *)
             | type '-w>' type   (* unrestricted, explicit *)
\end{verbatim}

\begin{verbatim}
Nat -> Nat                ;; simple function
Nat -> Nat -> Bool        ;; curried: Nat -> (Nat -> Bool)
[List A] -> Nat           ;; bracketed domain
[A -> B] -> [List A] -> List B  ;; higher-order
\end{verbatim}
\subsection{Dependent Types (Angle Brackets)}
\label{sec:orgb1acb83}

Angle brackets delimit dependent Pi (universal) and Sigma (existential) types:

\begin{verbatim}
dependent-type ::= '<' '(' ident ':' type ')' '->' type '>'  (* dependent Pi   *)
                 | '<' '(' ident ':' type ')' '*'  type '>'  (* dependent Sigma *)
                 | '<' ident ':' type '->' type '>'           (* shorthand Pi   *)
\end{verbatim}

\begin{verbatim}
;; Dependent function: "for all n : Nat, a vector of length n"
<(n : Nat) -> Vec Nat n>

;; Dependent pair: "a natural n together with a Vec of that length"
<(n : Nat) * Vec Nat n>

;; Shorthand (single binder, no inner parens):
<n : Nat -> Vec Nat n>
\end{verbatim}
\subsection{Sigma Types (Product / Pair Types)}
\label{sec:orgbdd2aa4}

Non-dependent pair types use infix \texttt{*}:

\begin{verbatim}
;; Non-dependent pair
Nat * Bool

;; In return type of split-at:
[Sigma [_ <List A>] [List A]]

;; Dependent Sigma via angle brackets:
<(x : Nat) * Vec Nat x>
\end{verbatim}
\subsection{Equality Type}
\label{sec:orgc865f24}

\begin{verbatim}
;; Eq type : proof that two terms are equal
(Eq Nat zero zero)       ;; 0 = 0 at type Nat
(Eq Nat [add x y] [add y x])  ;; commutativity
\end{verbatim}
\subsection{Union Types}
\label{sec:orga8b0290}

Union types use infix \texttt{|}:

\begin{verbatim}
Nat | Bool         ;; either a Nat or a Bool
Int | Rat | Nat    ;; right-associative: Int | (Rat | Nat)
\end{verbatim}
\subsection{Universe Levels}
\label{sec:org4dbafa4}

\begin{verbatim}
Type           ;; universe, level inferred
(Type 0)       ;; explicit level 0
(Type 1)       ;; explicit level 1 (contains Type 0)
\end{verbatim}
\subsection{Type Holes}
\label{sec:orge7bfb98}

The wildcard \texttt{\_} stands for an inferred type:

\begin{verbatim}
def x : _ zero    ;; type inferred as Nat
map _ xs          ;; type argument inferred
\end{verbatim}
\subsection{Typed Holes (Interactive Development)}
\label{sec:orgcc25878}

The \texttt{??} token is a \emph{typed hole} --- distinct from \texttt{\_} (inference hole). Where
\texttt{\_} means "infer this for me silently," \texttt{??} means "I don't know what goes here
--- show me what's possible." The type checker produces a diagnostic report with
the expected type, local context, and suggestions.

\begin{verbatim}
defn reverse
  | [nil]        -> nil
  | [[cons h t]] -> ??
  ;; Report: Hole ?? : [List A]
  ;;   h : A, t : [List A], reverse : [List A] -> [List A]
\end{verbatim}

Named holes label specific incomplete positions:

\begin{verbatim}
defn zip-with [f xs ys]
  match [xs ys]
    | [[cons x xs'] [cons y ys']] -> [cons ??combine ??recurse]
\end{verbatim}

\texttt{??} is the syntactic foundation for Editor-Assisted Interactive Development
(Agda/Idris-style hole-driven programming). See the Extended Spec Design
document for the full vision.
\section{Expressions}
\label{sec:org241c603}

\subsection{Function Application}
\label{sec:orgc4be96d}

Application is juxtaposition, with brackets \texttt{[]} for grouping:

\begin{verbatim}
;; Simple application
suc zero                    ;; suc(zero)

;; Multi-argument (uncurried convention)
add x y                     ;; add(x, y)

;; Brackets for grouping
[add x y]                   ;; same as add x y
cons [f a] [map f as]       ;; cons(f(a), map(f, as))

;; Nested application
map [add 1N _] '[1N 2N 3N]  ;; map with partial application
\end{verbatim}
\subsection{Lambda Expressions}
\label{sec:org5317ab4}

\begin{verbatim}
lambda ::= 'fn' binder+ body
         | '(' 'fn' '(' ident ':' type ')' body ')'   (* sexp mode *)
\end{verbatim}

\begin{verbatim}
;; WS mode: fn with bracket binder
fn [x : Nat] [suc x]

;; Multiple binder groups
fn [x : Nat] [y : Nat] [add x y]

;; Bare parameter (type inferred from context)
fn x [suc x]

;; Sexp mode: explicit binder
(fn (x : Nat) (suc x))

;; With multiplicity annotation
(fn (x :1 Nat) x)          ;; linear lambda
(fn (x :0 Nat) zero)       ;; erased lambda (can't use x)
\end{verbatim}
\subsection{Pattern Matching}
\label{sec:orgb870e5f}

\texttt{match} is the primary pattern matching form, using \texttt{|} for arm separation:

\begin{verbatim}
match-expr ::= 'match' scrutinee match-arm+
match-arm  ::= '|' pattern '->' body
\end{verbatim}

\begin{verbatim}
;; Match on Nat
match n
  | zero  -> zero
  | suc k -> add k k

;; Match on Bool
match b
  | true  -> suc zero
  | false -> zero

;; Match on List
match xs
  | nil       -> zero
  | cons a as -> suc [length as]

;; Match on Option
match [find pred xs]
  | none   -> default
  | some v -> v

;; Nested match
match xs
  | nil       -> nil
  | cons a as -> match [pred a]
    | true  -> cons a [filter pred as]
    | false -> filter pred as
\end{verbatim}

Available patterns:
\begin{itemize}
\item \textbf{Variable}: \texttt{x}, \texttt{n}, \texttt{rest} -- binds the matched value
\item \textbf{Wildcard}: \texttt{\_} -- matches anything, discards
\item \textbf{Constructor}: \texttt{zero}, \texttt{suc k}, \texttt{nil}, \texttt{cons a as}, \texttt{true}, \texttt{false},
\texttt{none}, \texttt{some v}, \texttt{ok v}, \texttt{err e}, \texttt{datum-nil}, \texttt{datum-cons h t}, etc.
\item \textbf{Nested}: \texttt{[cons x nil]} -- grouped pattern in brackets
\end{itemize}
\subsection{If (Conditional)}
\label{sec:org7d54272}

\texttt{if} is a macro that desugars to \texttt{match} on Bool:

\begin{verbatim}
;; Three-argument if
if [zero? n] base-case [step n]

;; Equivalent match:
match [zero? n]
  | true  -> base-case
  | false -> step n
\end{verbatim}
\subsection{Let (Local Binding)}
\label{sec:org93a96ed}

\texttt{let} is a macro that desugars to application:

\begin{verbatim}
;; Single binding
let x = [add 1N 2N]
  suc x

;; With type annotation
let x : Nat = [add 1N 2N]
  suc x
\end{verbatim}
\subsection{Pipe Operator (|>)}
\label{sec:org9a213fd}

The pipe operator threads a value left-to-right through functions:

\begin{verbatim}
;; Binary pipe
zero |> suc |> suc |> suc    ;; = suc(suc(suc(zero))) = 3

;; Block pipe with indented steps
|> nums
  map suc-fn
  filter is-positive
  reduce sum-rf zero

;; Block pipes support automatic loop fusion:
;; consecutive map/filter steps fuse into a single-pass
;; transducer when followed by a terminal (reduce, sum, etc.)
\end{verbatim}
\subsection{Compose Operator (>>)}
\label{sec:org139da18}

Left-to-right function composition:

\begin{verbatim}
;; Compose two functions
[suc >> suc] zero           ;; = suc(suc(zero)) = 2

;; Pipe into a composed function
zero |> [suc >> double]     ;; = double(suc(zero))
\end{verbatim}
\subsection{Type Annotations (the)}
\label{sec:org3f66276}

Explicit type annotation on an expression:

\begin{verbatim}
;; Annotate with explicit type
(the Nat zero)
(the [List Nat] nil)
\end{verbatim}
\subsection{Pairs}
\label{sec:org0cf7e33}

\begin{verbatim}
;; Construct a pair
pair 1N true

;; Project from a pair
first [pair 1N true]        ;; = 1N
second [pair 1N true]       ;; = true

;; Aliases: fst, snd
fst [pair x y]
snd [pair x y]
\end{verbatim}
\subsection{Quote and Quasiquote}
\label{sec:org5e347c7}

Quote \texttt{'} converts code to \texttt{Datum} (code-as-data):

\begin{verbatim}
'foo          ;; -> datum-sym
'42           ;; -> datum-nat
':hello       ;; -> datum-kw
'true         ;; -> datum-bool
'()           ;; -> datum-nil
'(add 1 2)   ;; -> datum-cons chain

;; Quasiquote with unquote
def val : Datum [datum-nat 10]
`,val          ;; splice val into template
`(add 1 ,val)  ;; template with hole
\end{verbatim}
\subsection{Collection Literals}
\label{sec:org9d91139}

\begin{verbatim}
;; List literal (linked list)
'[1N 2N 3N]

;; List with tail
'[1N 2N | rest-list]

;; PVec (persistent vector)
@[1 2 3 4 5]

;; Set
#{1 2 3}

;; Lazy sequence
~[1 2 3]

;; Map (key-value pairs)
{:name "Ada" :age 36}

;; Mixed-type map: value type auto-inferred as union
;; {:name "Alice" :age zero} : Map Keyword (Nat | String)
;; map-assoc widens: (map-assoc nat-map :k "s") : Map Keyword (Nat | String)
{:name "Alice" :age zero}
\end{verbatim}
\subsection{Partial Application}
\label{sec:orge3be61e}

Numbered placeholders \texttt{\_1}, \texttt{\_2} enable positional reordering:

\begin{verbatim}
;; Wildcard _ fills rightmost position
[add 1N _]           ;; fn x -> add 1N x

;; Numbered holes for reordering
[div _2 _1]          ;; fn x y -> div y x
\end{verbatim}
\subsection{Varargs}
\label{sec:org1eabae2}

The \texttt{...} marker in specs indicates variadic arguments. In \texttt{defn}, \texttt{...name}
collects excess arguments into a \texttt{List}:

\begin{verbatim}
spec sum-all Nat ... -> Nat
defn sum-all [...xs]
  sum xs

sum-all 1 2 3          ;; = 6
sum-all 1 2 3 4 5      ;; = 15

;; Mixed fixed + varargs
spec first-plus-rest Nat Nat ... -> Nat
defn first-plus-rest [first ...rest]
  add first [sum rest]

first-plus-rest 100 1 2 3  ;; = 106
\end{verbatim}
\section{Declarations}
\label{sec:org81b4991}

\subsection{Namespace (ns)}
\label{sec:org8c7a2b8}

Must be the first form in a file. Controls module identity and prelude loading:

\begin{verbatim}
;; Standard namespace: auto-imports prelude
ns my-project.utils

;; Bare namespace: no prelude
ns prologos.data.list :no-prelude
\end{verbatim}

The prelude automatically provides: \texttt{Nat}, \texttt{Bool}, \texttt{List}, \texttt{Option}, \texttt{Result},
\texttt{Pair} operations, \texttt{Eq=/=Ord=/=Add=/=Sub=/=Mul=/=Neg=/=Abs=/=FromInt=/=Num=/
=Fractional} traits and instances.
\subsection{Require}
\label{sec:org8ce16c0}

Import modules with named or aliased references:

\begin{verbatim}
;; Named imports
require [prologos.data.list :refer [List nil cons map length]]

;; Aliased module
require [prologos.data.nat :as nat :refer [add mult]]

;; Multiple requires
require [prologos.data.list   :refer [List nil cons map]]
        [prologos.data.nat    :refer [add mult]]
        [prologos.data.option :refer [Option none some]]
\end{verbatim}
\subsection{Value Definition (def)}
\label{sec:orge7fbea0}

\begin{verbatim}
;; With type annotation
def x : Nat zero

;; Shorthand with :=
def x := zero

;; Multi-line body (indented)
def nums : [List Nat]
  cons 1N [cons 2N [cons 3N [nil Nat]]]

;; Private (not auto-exported)
def- internal-state : Nat zero
\end{verbatim}
\subsection{Type Signature (spec)}
\label{sec:orgd8ae23a}

\texttt{spec} declares a type signature separately from the definition:

\begin{verbatim}
;; Simple function type
spec factorial Nat -> Nat

;; With implicit type parameters
spec map {A B : Type} [A -> B] [List A] -> List B

;; With trait constraints
spec elem {A : Type} [Eq A] A [List A] -> Bool

;; With varargs
spec sum-all Nat ... -> Nat

;; With multiple implicits and constraints
spec sort-on {A B : Type} [B -> B -> Bool] [A -> B] [List A] -> List A
\end{verbatim}
\subsubsection{Extended Spec with Metadata}
\label{sec:orgb5eb98f}

\texttt{spec} optionally accepts a trailing metadata map using the implicit map
syntax. Keyword-headed children after the type signature are collected into a
metadata hash. Recognized keys include \texttt{:implicits}, \texttt{:where}, \texttt{:doc},
\texttt{:examples}, \texttt{:properties}, \texttt{:see-also}, \texttt{:pre}, \texttt{:post}, \texttt{:refines}.
Unrecognized keys are stored but not acted upon (forward-compatible).

The \texttt{:implicits} key lifts implicit binders out of the type signature into
metadata, freeing the signature to express pure function shape. Inline
\texttt{\{A : Type\}} and \texttt{:implicits} coexist; both are merged.

\begin{verbatim}
;; With :implicits — clean signature:
spec sort [List A] -> [List A]
  :implicits {A : Type}
  :where (Ord A)
  :doc "Sorts a list in ascending order"
  :examples
    - [sort '[3N 1N 2N]] => '[1N 2N 3N]
    - [sort '[]] => '[]
  :properties (sortable-laws A)
  :see-also [reverse filter]

;; HKT case — :implicits dramatically improves readability:
spec gmap [A -> B] -> [C A] -> [C B]
  :implicits {A B : Type} {C : Type -> Type}
  :where (Seqable C) (Buildable C)

;; With functor — Pi types hidden entirely:
spec xf-compose [Xf A B] -> [Xf B C] -> [Xf A C]
  :implicits {A B C : Type}
  :properties (transducer-fusion-laws A B C)
\end{verbatim}

The metadata is entirely optional. A simple \texttt{spec add Nat Nat -> Nat} is
unchanged. Metadata keys follow the principle of progressive disclosure: add
\texttt{:doc} and \texttt{:examples} for public APIs, \texttt{:properties} for libraries,
\texttt{:pre=/}:post= for safety-critical code, \texttt{:refines} for verified code.

Each \texttt{:properties} entry is a map with three keys:

\begin{center}
\begin{tabular}{ll}
Key & Meaning\\
\hline
\texttt{:name} & Human-readable label for the property\\
\texttt{:forall} & Universally quantified variables (binders)\\
\texttt{:holds} & Boolean expression that must be true\\
\end{tabular}
\end{center}

See the Extended Spec Design document
(\texttt{docs/tracking/2026-02-22\_EXTENDED\_SPEC\_DESIGN.org}) for the full research
survey and design rationale.
\subsection{Function Definition (defn)}
\label{sec:orgfdd433b}

\begin{verbatim}
;; Basic function
spec add Nat -> Nat -> Nat
defn add [x y]
  match x
    | zero  -> y
    | suc n -> suc [add n y]

;; With inline type
defn suc-all [xs : List Nat] : List Nat
  map [add _ 1N] xs

;; With return type annotation
impl Eq Nat
  defn eq? [x y] <Bool>
    nat-eq? x y

;; Multi-arity with |
defn fact
  | [zero]  -> suc zero
  | [suc n] -> mult [suc n] [fact n]

;; Private
defn- helper [x]
  suc x
\end{verbatim}
\subsection{Algebraic Data Types (data)}
\label{sec:orgc04f484}

\begin{verbatim}
;; Enumeration (nullary constructors)
data Bool
  true
  false

;; Parameterized with constructors
data List {A}
  nil
  cons : A -> List A

;; Multiple implicit params
data Either {A B}
  left  : A
  right : B

;; Single constructor (wrapper)
data Ordering
  lt-ord
  eq-ord
  gt-ord

;; Private data type
data- InternalTree {A}
  leaf
  node : A -> InternalTree A -> InternalTree A
\end{verbatim}
\subsection{Traits and Implementations}
\label{sec:orgfcc3bfe}

\begin{verbatim}
;; Trait definition
trait Eq {A}
  eq? : A -> A -> Bool

;; Implementation
impl Eq Nat
  defn eq? [x y] <Bool>
    nat-eq? x y

;; Multi-method trait
trait Ord {A}
  compare : A -> A -> Ordering

;; Generic functions using trait dicts
spec eq-neq [Eq A] A A -> Bool
defn eq-neq [dict x y]
  not [dict x y]

;; Bundle: named trait combination
bundle Num := (Add, Sub, Mul, Neg, Abs, FromInt)

;; Trait with :laws reference
trait Functor {F : Type -> Type}
  :laws (functor-laws F)
  fmap : {A B : Type} [A -> B] -> [F A] -> [F B]
\end{verbatim}
\subsection{Property Declarations}
\label{sec:org61c8206}

A \texttt{property} declares a named, composable group of propositions --- analogous
to \texttt{bundle} for traits. Properties compose via \texttt{:includes} (conjunction) and
attach to specs via \texttt{:properties} or to traits via \texttt{:laws}.

\begin{verbatim}
;; Basic property declaration
property sortable-laws {A : Type}
  :where (Ord A)
  - :name "idempotent"
    :forall {xs : [List A]}
    :holds [eq? [sort [sort xs]] [sort xs]]
  - :name "length-preserving"
    :forall {xs : [List A]}
    :holds [eq? [length [sort xs]] [length xs]]

;; Composition via :includes (like bundle for traits)
property monoid-laws {A : Type}
  :where (Add A) (AdditiveIdentity A)
  :includes (semigroup-laws A)
  - :name "left-identity"
    :forall {x : A}
    :holds [eq? [add additive-identity x] x]
  - :name "right-identity"
    :forall {x : A}
    :holds [eq? [add x additive-identity] x]

;; Hierarchical: monad includes applicative includes functor
property monad-laws {M : Type -> Type}
  :where (Monad M)
  :includes (applicative-laws M)
  - :name "left-identity"
    :forall {a : A} {f : [A -> [M B]]}
    :holds [eq? [bind [pure a] f] [f a]]

;; Usage in spec
spec sort [List A] -> [List A]
  :implicits {A : Type}
  :where (Ord A)
  :properties (sortable-laws A)

;; Higher-order: parameterized over functions
property sorting-properties {A : Type} {f : [List A] -> [List A]}
  :where (Ord A)
  - :name "idempotent"
    :forall {xs : [List A]}
    :holds [eq? [f [f xs]] [f xs]]

spec sort [List A] -> [List A]
  :where (Ord A)
  :properties (sorting-properties A sort)
\end{verbatim}

Property clause names are scoped to their declaring block. When \texttt{:includes}
flattens a hierarchy, names are qualified with \texttt{/}:
\texttt{functor-laws/identity}, \texttt{monad-laws/left-identity}, etc.
\subsection{Functor Declarations (Named Type Abstractions)}
\label{sec:org2c36132}

A \texttt{functor} declares a named, parameterized type abstraction with optional
category-theoretic metadata. Transparent by default --- the type checker unfolds
to the \texttt{:unfolds} form during elaboration. Metadata keys map CT concepts to
approachable language.

\begin{verbatim}
;; Simple type synonym
functor FilePath
  :unfolds String

;; Parameterized synonym with documentation
functor Result {A}
  :unfolds [Either String A]
  :doc "A computation that may fail with a string error"

;; Full category-theoretic treatment
functor Xf {A B : Type}
  :doc "A transducer: transforms A-reductions into B-reductions"
  :compose xf-compose
  :identity id-xf
  :laws (transducer-fusion-laws A B)
  :unfolds <(S :0 Type) -> [S -> B -> S] -> S -> A -> S>

;; Optics
functor Lens {S T A B : Type}
  :doc "A bidirectional accessor: view and update a part of a structure"
  :compose lens-compose
  :identity lens-id
  :laws (lens-laws S T A B)
  :see-also [Prism Traversal Iso]
  :unfolds <{F : Type -> Type} -> (Functor F) -> [A -> [F B]] -> S -> [F T]>

;; Usage in specs — the functor name replaces raw Pi types:
spec xf-compose [Xf A B] -> [Xf B C] -> [Xf A C]
  :implicits {A B C : Type}
spec into-list [Xf A B] -> [List A] -> [List B]
  :implicits {A B : Type}
\end{verbatim}

\begin{center}
\begin{tabular}{lll}
CT concept & Key & Meaning in plain English\\
\hline
Object mapping & \texttt{:unfolds} & "What this expands to under the hood"\\
Morphism composition & \texttt{:compose} & "How to chain two of these together"\\
Identity morphism & \texttt{:identity} & "The do-nothing version"\\
Laws & \texttt{:laws} & "What rules these always follow"\\
\end{tabular}
\end{center}

Only \texttt{:unfolds} is required. The rest are progressive.
\subsection{User-Defined Macros}
\label{sec:orgc1f4c3a}

\begin{verbatim}
;; Simple macro: double-it x --> add x x
defmacro double-it [$x] [add $x $x]

double-it 5             ;; = 10

;; Conditional macro
defmacro when-nonzero [$n $body]
  [if [not [zero? $n]] $body zero]

when-nonzero 5 [double 5]    ;; = 10
when-nonzero zero [double 5]  ;; = 0

;; Introspection
(expand (double-it 5))        ;; shows expansion
(expand-1 (double-it 5))      ;; one step
(expand-full (when true [suc zero]))  ;; all steps with labels
\end{verbatim}
\section{Multiplicity (QTT)}
\label{sec:org2147f01}

Prologos uses Quantitative Type Theory for resource tracking. Each variable
binding has a \emph{multiplicity} annotation:

\begin{center}
\begin{tabular}{lll}
Annotation & Meaning & Syntax\\
\hline
\texttt{:0} & Erased (0 uses at runtime) & \texttt{(x :0 Nat)}\\
\texttt{:1} & Linear (exactly 1 use) & \texttt{(x :1 Nat)}\\
\texttt{:w} & Unrestricted (any uses) & \texttt{(x :w Nat)}\\
\end{tabular}
\end{center}

Arrow types carry multiplicities:

\begin{verbatim}
Nat -> Nat            ;; unrestricted (default :w)
Nat -0> Nat           ;; erased: argument not used at runtime
Nat -1> Nat           ;; linear: argument used exactly once
Nat -w> Nat           ;; unrestricted: explicit
\end{verbatim}

\begin{verbatim}
;; Erased argument (type-level only)
(fn (A :0 Type) (fn (x :w A) x))     ;; polymorphic identity

;; Linear function
(fn (x :1 Nat) x)                    ;; must use x exactly once
\end{verbatim}
\section{Dependent Types and Eliminators}
\label{sec:org3677728}

\subsection{Natural Number Elimination (natrec)}
\label{sec:org57d8141}

The low-level eliminator for natural numbers:

\begin{verbatim}
;; natrec motive base step target
;; motive : Nat -> Type
;; base   : motive zero
;; step   : (n : Nat) -> motive n -> motive (suc n)
;; target : Nat
natrec motive base step target
\end{verbatim}

In practice, \texttt{match} is preferred:

\begin{verbatim}
;; Idiomatic: match instead of natrec
defn double [n]
  match n
    | zero  -> zero
    | suc k -> suc [suc [double k]]
\end{verbatim}
\subsection{Equality Elimination (J)}
\label{sec:org2d7a57c}

\begin{verbatim}
;; J motive base left right proof
;; Given proof : Eq A left right
;; J computes: motive(left, right, proof) via base(left)
J motive base left right refl
\end{verbatim}
\subsection{Length-Indexed Vectors}
\label{sec:org5a7718c}

\begin{verbatim}
;; Vec A n : a vector of exactly n elements of type A
;; Fin n   : a number guaranteed less than n

vnil Nat                         ;; empty: Vec Nat 0
vcons Nat 2N x xs                ;; cons:  Vec Nat 3 (given xs : Vec Nat 2)
vhead Nat 2N v                   ;; safe head (non-empty guaranteed)
vtail Nat 2N v                   ;; safe tail
vindex Nat 3N [fzero 2N] v      ;; safe index via Fin
\end{verbatim}
\section{Full Program Example}
\label{sec:orgf02f979}

A complete Prologos program demonstrating multiple features:

\begin{verbatim}
ns examples.demo

require [prologos.data.list   :refer [List nil cons map reduce sum length]]
        [prologos.data.nat    :refer [add mult]]
        [prologos.data.option :refer [Option none some]]

;; Type signature with implicit params
spec factorial Nat -> Nat
defn factorial [n]
  match n
    | zero  -> suc zero
    | suc k -> mult n [factorial k]

;; Higher-order function with trait constraint
spec sum-mapped {A : Type} [A -> Nat] [List A] -> Nat
defn sum-mapped [f xs]
  sum [map f xs]

;; Data type definition
data Tree {A}
  leaf
  node : A -> Tree A -> Tree A

;; Pattern matching on user-defined types
spec tree-size [Tree A] -> Nat
defn tree-size [t]
  match t
    | leaf       -> zero
    | node _ l r -> suc [add [tree-size l] [tree-size r]]

;; Using pipe operator
def result : Nat
  |> '[1N 2N 3N 4N 5N]
    map [add 1N _]
    reduce add 0N

;; Macro definition
defmacro unless [$cond $body] [if $cond zero $body]

;; Using the macro
unless [zero? 5] [factorial 5]
\end{verbatim}
\section{Appendix: S-Expression Mode}
\label{sec:org80043a0}

Every WS-mode form has a canonical s-expression representation. In sexp mode,
all grouping uses parentheses and whitespace is not significant:

\begin{verbatim}
;; WS mode:
defn add [x y]
  match x
    | zero  -> y
    | suc n -> suc [add n y]

;; Equivalent sexp mode:
(defn add (x y)
  (match x
    (zero  -> y)
    (suc n -> (suc (add n y)))))

;; WS mode:
spec map {A B : Type} [A -> B] [List A] -> List B

;; Equivalent sexp mode:
(spec map {A B : Type} (-> (-> A B) (-> (List A) (List B))))
\end{verbatim}
\section{Appendix: Whitespace Reader Rules}
\label{sec:orgb8dc05a}

The WS reader converts indentation to explicit structure:

\begin{enumerate}
\item \textbf{Same indentation} as previous line: NEWLINE token (sibling separator)
\item \textbf{Deeper indentation}: INDENT token (child block begins)
\item \textbf{Shallower indentation}: DEDENT token(s) (child blocks end)
\item \textbf{Inside brackets} \texttt{[]}, \texttt{()}, \texttt{\{\}}, \texttt{<>}: newlines are whitespace
\item \textbf{Blank/comment lines}: ignored for indentation purposes
\item \textbf{Column 0}: always starts a new top-level form
\end{enumerate}
\section{Appendix: Reader Desugaring Table}
\label{sec:orgd158814}

\begin{center}
\begin{tabular}{lll}
Surface syntax & Reader output & Meaning\\
\hline
\texttt{'expr} & \texttt{(\$quote expr)} & Quote to Datum\\
\texttt{`expr} & \texttt{(\$quasiquote expr)} & Quasiquote template\\
\texttt{,expr} & \texttt{(\$unquote expr)} & Unquote (splice into QQ)\\
\texttt{42N} & \texttt{(\$nat-literal 42)} & Natural number literal\\
\texttt{\textasciitilde{}42} & \texttt{(\$approx-literal 42)} & Approximate (Posit) literal\\
\texttt{3.14} & \texttt{(\$decimal-literal 157/50)} & Bare decimal → Posit32\\
\texttt{:name} & \texttt{keyword token} & Keyword literal\\
\texttt{'[1 2 3]} & \texttt{(\$list-literal 1 2 3)} & List literal\\
\texttt{@[1 2 3]} & \texttt{(\$vec-literal 1 2 3)} & PVec literal\\
\texttt{\textasciitilde{}[1 2 3]} & \texttt{(\$lseq-literal 1 2 3)} & Lazy sequence literal\\
\texttt{\#\{1 2 3\}} & \texttt{(\$set-literal 1 2 3)} & Set literal\\
\texttt{...name} & \texttt{(\$rest-param name)} & Rest/vararg parameter\\
\texttt{x \textbackslash{}vert> f} & \texttt{(\$pipe-gt x f)} & Pipe left-to-right\\
\texttt{f >{}>{} g} & \texttt{(\$compose f g)} & Compose left-to-right\\
\end{tabular}
\end{center}
\end{document}
